\subsection{Setup}\label{sec:setup}

\paragraph{Data.}
\emph{Coat}~\cite{schnabel2016} has 290 users, 300 items and 6{,}960 self-selected
(MNAR) ratings, exactly 24 per user. The test data are 16 ratings per user on
uniformly random items (MAR); we rank only the items the user did not rate,
which leaves 14.7 candidates per user on average, and 225 users have at least
one positive candidate. The dataset ships propensities, which we use.
\emph{Yahoo!\,R3}~\cite{marlin2009collaborative} has 311{,}704 self-selected ratings of
15{,}400 users on 1{,}000 songs (2.0\% dense) and 54{,}000 ratings of 10 random
songs for 5{,}400 of the users, in the version distributed
with~\cite{chen2021autodebias}. A random 5\% of the MAR users is held out and used
only to fit an outcome-dependent Naive-Bayes propensity~\cite{schnabel2016}; in a
preliminary comparison on validation data it did not improve the propagation
operators, and all reported results use the logistic exposure model. The remaining 5{,}130 test users have 10 candidates each, and 2{,}315 of
them have a positive. \emph{KuaiRec}~\cite{gao2022kuairec} provides the big matrix
(7{,}176 users, 10{,}728 items, 10.3M logged views, 13.4\% dense) as the MNAR log and
a fully observed small matrix (1{,}411 users, 3{,}327 items) as the test set. We rank
all small-matrix items that the user never viewed in the log (3{,}314 candidates
per user, 154 of them positive on average). \emph{KuaiRand-Pure}~\cite{gao2022kuairand}
is the only one of the four datasets whose log records platform exposure rather
than self-selection: $O$ is an impression of the standard recommendation feed and
$Y$ a click. We keep a random 10{,}000 of the users who also received randomly
exposed videos, and all 7{,}583 videos; the standard log then has 511{,}863
impressions (0.68\% dense, click rate 47\%). The random impressions of the same
period, restricted to videos the user was not shown in the standard log, are
the test set (43.6 candidates per user, 17.5\% clicked; 8{,}347 users have a
positive). Positives are ratings $\ge4$ (Coat, Yahoo!\,R3), a watch ratio $\ge2$
(KuaiRec) and a click (KuaiRand). A random ranking reaches nDCG@5 of 0.306 on
Coat and 0.352 on Yahoo!\,R3 and nDCG@20 of 0.046 on KuaiRec. The candidate sets
of Coat and Yahoo!\,R3 are small, so their metric differences are compressed;
KuaiRec ranks thousands of candidates. We report nDCG@5 on Coat and
Yahoo!\,R3, nDCG@10 on KuaiRand and nDCG@20 on KuaiRec, and AUC and Recall in
the Supplementary Material. On Coat and Yahoo!\,R3 the log is self-selected
ratings: ``exposure'' there means that the user chose to rate the item, and
exposure invariance is invariance to that choice, not to a platform's
allocation.

\paragraph{Protocol.}
The unbiased data are split into a validation part (30\% of each user's
candidates on Coat, 30\% of the test users otherwise) and a test part, over 10
(Coat) or 5 random splits. Every method selects its configuration on the
validation part of each split, and trained models also select their epoch
there: each split stops on its own validation score and never sees the
validation or test users of another split. All fitted nuisances are cross-fitted over ten folds of pairs, except
for the sample-split variants, whose imputation (and propensity model, where
it is estimated) is fitted on a random 20\% of the pairs; only the latter satisfy the assumptions of the theory
(Remark~\ref{rem:xfit}). We report the mean and standard deviation over splits
of the test metric. Because paired tests over re-sampled splits are anti-conservative when the
splits overlap~\cite{nadeau2003}, differences are tested at the level
of users, with a Holm correction over all comparisons of a dataset
(Appendix~\ref{app:details}). These tests treat the training log, the fitted
nuisances and the trained models as fixed: they account for the variation
over test users, not for re-drawing the log or re-fitting the pipeline, which
we assess only through the split and seed standard deviations. Progress claims in top-$N$ recommendation have
repeatedly failed against well-tuned simple
baselines~\cite{dacrema2019progress,rendle2019baselines,anelli2023myth}; we therefore
tune every method on the same validation data and report how many
configurations each one searched.

\paragraph{Methods.}
We compare DRUP with up to 26 baselines (22 on KuaiRand, where eleven models were trained). \emph{Non-personalised}: popularity and the imputation
$\hat Y$ alone. \emph{Trained with a pointwise loss}: MF, IPS-MF~\cite{schnabel2016},
DR-MF, DR-JL~\cite{wang2019drjl}, MRDR~\cite{guo2021mrdr}, iALS~\cite{hu2008ials},
LightGCN and DR-LightGCN (LightGCN with the DR loss).
\emph{Trained with a pairwise or contrastive loss}: MF, PDA~\cite{zhang2021pda},
MACR~\cite{wei2021macr}, LightGCN~\cite{he2020lightgcn}, r-AdjNorm~\cite{zhao2022radjnorm},
NAVIP~\cite{kim2022navip} and SimGCL~\cite{yu2022simgcl}, all with BPR.
\emph{Training-free on the logged graph}: linear LightGCN
(the one- plus three-hop form of Section~\ref{sec:drup} on $O\odot Y$; with the searched
normalisation exponent it is a linear r-AdjNorm), EASE~\cite{steck2019ease},
GF-CF~\cite{shen2021gfcf} and BSPM~\cite{choi2023bspm}. \emph{Training-free on a debiased graph}: IPS adjacency
(NAVIP-style), IPS adjacency with the walk correction, EASE, GF-CF and BSPM on the DR
graph $W$, the uncorrected DR adjacency, and the five-hop DR adjacency. Besides
DRUP we report its five-hop version and DR adjacency and DRUP with the sample
split of Section~\ref{sec:drup} (DRUP-split), over five draws of the split-off
pairs. Grids, the number of configurations per method, and training details are
in Appendix~\ref{app:details}. The training-free methods search between 1 and
6{,}300 configurations and the trained ones between 6 and 108. Trained models
are evaluated after every epoch; each split keeps the best (configuration,
epoch) on its own validation part and stops once its validation score has not
improved for a fixed number of epochs.
